\section{Zadanie 4}
Niech \(A_0, A_1, \dots\) będzie nieskończonym ciągiem zdarzeń i
niech \(P(A_n) = 1\) dla każdego \(n\).
Udowodnij, że \(P(\bigcap_{n=0}^{\infty}A_n) = 1\)

\subsection*{Rozwiązanie}
Dla każdego n mamy \(P(A_n) = 1\), więc dla każdego dopełnienia zdarzenia \(A'_n = \Omega - A_n\) mamy:
\(P(A'_n) = 1 - P(A_n) = 0\).

Zauważmy, że dopełnienie przecięcia jest równe sumie dopełnień
\[
	(\bigcap_{n=0}^{\infty}A_n)' = \bigcup_{n=0}^{\infty}A_n'
\]
Teraz korzystamy z nierówności pomiędzy prawdopodobieństwem sumy a sumą prawdopodobieństw (Zadanie 3),
z czego ładnie nam wychodzi teza
\[
	P(\bigcup_{n=0}^{\infty}A_n') \leq \Sigma_{n=0}^{\infty}P(A'_n) = \Sigma_{n=0}^{\infty}0 = 0
\]
\[
	P(\bigcup_{n=0}^{\infty}A_n') = 0
\]
\[
	P(\bigcap_{n=0}^{\infty}A_n) = 1 - 0 = 1
\]
